% ---
% title: Two identical dice
% tags: stars-and-bars, combinations
% ---

\begin{problem}
	We roll two identical (indistinguishable) dice. How many different outcomes
	can we observe?
\end{problem}

\begin{slides}
	\small
	\exstatement
	\pause
	\sol
	$(1,2)$ and $(2,1)$ look the same. Pairs with two different numbers:
	$\binom62=15$; doubles: $6$; total $21$.\\[1mm]
	Or: choose 2 out of 6 numbers, order irrelevant, repetition allowed (stars
	and bars, as in Exercise~\exref{stars-and-bars}):
	$\binom{6+2-1}{2}=\binom72 = 21$.
\end{slides}

\begin{solution}
	\idea Because the dice look the same, only \emph{which} numbers appear
	matters, not on which die.

	\step{First way.} $(1,2)$ and $(2,1)$ are one observation. Outcomes with two
	different numbers: choose 2 of 6 numbers, $\binom62=15$. Doubles: $6$.
	Total $15+6=21$.

	\step{Second way.} We choose 2 numbers out of 6, order irrelevant,
	repetition allowed. This is stars and bars (Exercise~\exref{stars-and-bars}):
	2 stars distributed among the 6 numbers, $\binom{6+2-1}{2}=\binom72=21$.
	\begin{answer}
		$21$
	\end{answer}
	\begin{remark}
		\textbf{Common mistake:} $36/2=18$. Halving is wrong because the doubles
		$(1,1),\dots,(6,6)$ appear only once among the 36 ordered outcomes, not
		twice: $(36-6)/2+6=21$.
	\end{remark}
\end{solution}
